% GENERATED FILE. Edit canonical lessons, then run npm run generate:books.
% canonical-source-hash: fnv1a64:120d819c063b8924
% canonical-lessons: BN-C282-churi-kora, BN-C282-goyna, BN-C282-pakrano, BN-C282-thela, BN-C282-poketmar

\chapter{A Theft in the Night}
\label{ch:bn-a-theft-in-the-night}
\hlchaptermodality{282}

\begin{chapteropening}
\textbf{By the end of this chapter:} \emph{I can tell how a thief got in and what was taken.}

Complete the last lesson of chapter 282: I can tell how a thief got in and what was taken.
\end{chapteropening}

\section[\texorpdfstring{\bn{চুরি} \bn{করা} (churi kôrā)}{চুরি করা (churi kôrā)}]{\bn{চুরি} \bn{করা} (churi kôrā) — to steal}

\label{lesson:BN-C282-churi-kora}



\begin{quote}
Before the new one: say the Bengali for to rush, then the Bengali for at last.
\end{quote}

\subsection*{You'll want to know: \bn{চুরি} \bn{করা}}
\textbf{\bn{চুরি} \bn{করা}} — \emph{churi kôrā} — \textquotedblleft{}to steal\textquotedblright{}.

\begin{quote}
\bn{দরজায়} \bn{তালা} \bn{দিতে} \bn{হবে}, \bn{নাহলে} \bn{চোর} \bn{ঢুকে} \bn{সব} \bn{চুরি} \bn{করবে।}
\end{quote}

\begin{quote}
\emph{dôrjāy tālā dite hôbe, nāhôle chor ḍhuke sôb churi kôrbe.}
\end{quote}

\begin{quote}
We have to lock the door, otherwise a thief will get in and steal everything.
\end{quote}

Past (he/she …): \bn{সেই} \bn{চোর} \bn{টাকা} \bn{চুরি} \bn{করল}, \emph{sei chor ṭākā churi kôrlô}.

\begin{cousinweb}
\textbf{\bn{চুরি} \bn{করা}} (\emph{churi kôrā}) puts \bn{চুরি}, 'a theft' — Sanskrit \emph{caurikā} (CDIAL 4937, which lists Bengali \emph{curi}), the same family as \textbf{\bn{চোর}}, a thief — before \bn{করা}: to do a theft.
\end{cousinweb}

\begin{grammarlens}[title={What you've built}]
One more word for telling what happened.
\end{grammarlens}

\subsection*{Your turn}
\begin{itemize}
\raggedright
  \item \emph{Say it:} \emph{churi kôrā}
  \item \emph{Say it:} \emph{churi kôrā}, once more
  \item \emph{Say it:} \emph{dôrjāy tālā dite hôbe, nāhôle chor ḍhuke sôb churi kôrbe.} — We have to lock the door, otherwise a thief will get in and steal everything.
  \item \emph{Recall:} say the Bengali for to rush, then the Bengali for at last, then say \emph{churi kôrā} again
\end{itemize}

\begin{tcolorbox}[breakable,colback=teal!4,colframe=teal!35!black,title={Before you move on}]
What does \bn{চুরি} \bn{করা} mean? (\textquotedblleft{}To steal\textquotedblright{}.) Say it once more.
\end{tcolorbox}
\section[\texorpdfstring{\bn{গয়না} (gôynā)}{গয়না (gôynā)}]{\bn{গয়না} (gôynā) — jewellery; an ornament}

\label{lesson:BN-C282-goyna}



\begin{quote}
Before the new one: say the Bengali for at last, then the Bengali for to steal.
\end{quote}

\subsection*{You'll want to know: \bn{গয়না}}
\textbf{\bn{গয়না}} — \emph{gôynā} — \textquotedblleft{}jewellery; an ornament\textquotedblright{}.

\begin{quote}
\bn{চোর} \bn{আলমারি} \bn{খুলে} \bn{দিদিমার} \bn{গয়না} \bn{নিয়ে} \bn{পালাল।}
\end{quote}

\begin{quote}
\emph{chor ālmāri khule didimār gôynā niye pālālô.}
\end{quote}

\begin{quote}
The thief opened the cupboard and ran off with my grandmother's jewellery.
\end{quote}

\begin{cousinweb}
\textbf{\bn{গয়না}} (\emph{gôynā}) is filed in CDIAL 4364 under Sanskrit \emph{grahaṇa}, 'seizing, holding; putting on clothes', where Bengali \emph{gaynā} appears; how it came to mean an ornament is not set out there.
\end{cousinweb}

\begin{grammarlens}[title={What you've built}]
One more word for telling what happened.
\end{grammarlens}

\subsection*{Your turn}
\begin{itemize}
\raggedright
  \item \emph{Say it:} \emph{gôynā}
  \item \emph{Say it:} \emph{gôynā}, once more
  \item \emph{Say it:} \emph{chor ālmāri khule didimār gôynā niye pālālô.} — The thief opened the cupboard and ran off with my grandmother's jewellery.
  \item \emph{Recall:} say the Bengali for at last, then the Bengali for to steal, then say \emph{gôynā} again
\end{itemize}

\begin{tcolorbox}[breakable,colback=teal!4,colframe=teal!35!black,title={Before you move on}]
What does \bn{গয়না} mean? (\textquotedblleft{}Jewellery; an ornament\textquotedblright{}.) Say it once more.
\end{tcolorbox}
\section[\texorpdfstring{\bn{পাকড়ানো} (pākṛāno)}{পাকড়ানো (pākṛāno)}]{\bn{পাকড়ানো} (pākṛāno) — to grab; to nab}

\label{lesson:BN-C282-pakrano}



\begin{quote}
Before the new one: say the Bengali for to steal, then the Bengali for jewellery.
\end{quote}

\subsection*{You'll want to know: \bn{পাকড়ানো}}
\textbf{\bn{পাকড়ানো}} — \emph{pākṛāno} — \textquotedblleft{}to grab; to nab\textquotedblright{}.

\begin{quote}
\bn{প্রতিবেশীরা} \bn{চোরকে} \bn{রাস্তায়} \bn{পাকড়াল।}
\end{quote}

\begin{quote}
\emph{prôtibeshīrā chorke rāstāy pākṛālô.}
\end{quote}

\begin{quote}
The neighbours grabbed the thief in the street.
\end{quote}

Past (he/she …): \bn{সেই} \bn{লোক} \bn{চোরকে} \bn{পাকড়াল}, \emph{sei lok chorke pākṛālô}.

\begin{cousinweb}
\textbf{\bn{পাকড়ানো}} (\emph{pākṛāno}) is filed in CDIAL 7619 under a Middle Indo-Aryan \emph{pakkaḍ-}, 'seize', which lists Bengali \emph{pākaṛā}, 'to seize'. It is everyday, rather rough talk for catching hold of someone.
\end{cousinweb}

\begin{grammarlens}[title={What you've built}]
One more word for telling what happened.
\end{grammarlens}

\subsection*{Your turn}
\begin{itemize}
\raggedright
  \item \emph{Say it:} \emph{pākṛāno}
  \item \emph{Say it:} \emph{pākṛāno}, once more
  \item \emph{Say it:} \emph{prôtibeshīrā chorke rāstāy pākṛālô.} — The neighbours grabbed the thief in the street.
  \item \emph{Recall:} say the Bengali for to steal, then the Bengali for jewellery, then say \emph{pākṛāno} again
\end{itemize}

\begin{tcolorbox}[breakable,colback=teal!4,colframe=teal!35!black,title={Before you move on}]
What does \bn{পাকড়ানো} mean? (\textquotedblleft{}To grab; to nab\textquotedblright{}.) Say it once more.
\end{tcolorbox}
\section[\texorpdfstring{\bn{ঠেলা} (ṭhelā)}{ঠেলা (ṭhelā)}]{\bn{ঠেলা} (ṭhelā) — to push; to shove}

\label{lesson:BN-C282-thela}



\begin{quote}
Before the new one: say the Bengali for jewellery, then the Bengali for to grab.
\end{quote}

\subsection*{You'll want to know: \bn{ঠেলা}}
\textbf{\bn{ঠেলা}} — \emph{ṭhelā} — \textquotedblleft{}to push; to shove\textquotedblright{}.

\begin{quote}
\bn{চোর} \bn{আমাকে} \bn{ঠেলে} \bn{দরজা} \bn{দিয়ে} \bn{পালাল।}
\end{quote}

\begin{quote}
\emph{chor āmāke ṭhele dôrjā diye pālālô.}
\end{quote}

\begin{quote}
The thief shoved me aside and escaped through the door.
\end{quote}

Past (he/she …): \bn{সেই} \bn{চোর} \bn{আমাকে} \bn{ঠেলল}, \emph{sei chor āmāke ṭhellô}.

\begin{cousinweb}
\textbf{\bn{ঠেলা}} (\emph{ṭhelā}) is filed in CDIAL 5512 under a Middle Indo-Aryan \emph{ṭhēḍḍ-}, 'push', which lists Bengali \emph{ṭhelā}, 'to push'. A \bn{ঠেলাগাড়ি} is a handcart, a cart you push.
\end{cousinweb}

\begin{grammarlens}[title={What you've built}]
One more word for telling what happened.
\end{grammarlens}

\subsection*{Your turn}
\begin{itemize}
\raggedright
  \item \emph{Say it:} \emph{ṭhelā}
  \item \emph{Say it:} \emph{ṭhelā}, once more
  \item \emph{Say it:} \emph{chor āmāke ṭhele dôrjā diye pālālô.} — The thief shoved me aside and escaped through the door.
  \item \emph{Recall:} say the Bengali for jewellery, then the Bengali for to grab, then say \emph{ṭhelā} again
\end{itemize}

\begin{tcolorbox}[breakable,colback=teal!4,colframe=teal!35!black,title={Before you move on}]
What does \bn{ঠেলা} mean? (\textquotedblleft{}To push; to shove\textquotedblright{}.) Say it once more.
\end{tcolorbox}
\section[\texorpdfstring{\bn{পকেটমার} (pôkeṭmār)}{পকেটমার (pôkeṭmār)}]{\bn{পকেটমার} (pôkeṭmār) — a pickpocket}

\label{lesson:BN-C282-poketmar}



\begin{quote}
Before the new one: say the Bengali for to grab, then the Bengali for to push.
\end{quote}

\subsection*{You'll want to know: \bn{পকেটমার}}
\textbf{\bn{পকেটমার}} — \emph{pôkeṭmār} — \textquotedblleft{}a pickpocket\textquotedblright{}.

\begin{quote}
\bn{বাসের} \bn{ভিড়ে} \bn{কোনো} \bn{পকেটমার} \bn{আমার} \bn{মানিব্যাগ} \bn{নিল।}
\end{quote}

\begin{quote}
\emph{bāser bhiṛe kono pôkeṭmār āmār mānibyāg nilô.}
\end{quote}

\begin{quote}
In the crowded bus some pickpocket took my wallet.
\end{quote}

\begin{cousinweb}
\textbf{\bn{পকেটমার}} (\emph{pôkeṭmār}) joins \textbf{\bn{পকেট}}, a pocket, from English, to \bn{মার}, the stem of \textbf{\bn{মারা}}, which in slang also means to pinch: one who hits pockets.
\end{cousinweb}

\begin{grammarlens}[title={What you've built}]
One more word for telling what happened.
\end{grammarlens}

\subsection*{Your turn}
\begin{itemize}
\raggedright
  \item \emph{Say it:} \emph{pôkeṭmār}
  \item \emph{Say it:} \emph{pôkeṭmār}, once more
  \item \emph{Say it:} \emph{bāser bhiṛe kono pôkeṭmār āmār mānibyāg nilô.} — In the crowded bus some pickpocket took my wallet.
  \item \emph{Recall:} say the Bengali for to grab, then the Bengali for to push, then say \emph{pôkeṭmār} again
\end{itemize}

\begin{tcolorbox}[breakable,colback=teal!4,colframe=teal!35!black,title={Before you move on}]
What does \bn{পকেটমার} mean? (\textquotedblleft{}A pickpocket\textquotedblright{}.) Say it once more.
\end{tcolorbox}
